\documentclass[10pt]{article}

\def\Type{LessonDJ}
\def\TypeNum{Workshop Week 12} %Lesson #
\def\TypeTitle{Probability and Optimization} %Lesson Title
\def\Semester{Fall 2021} %Not Used 
\def\Course{MATH 1190}%Course number and name
\def\Targets{    \target{I2}, \target{P1}, \target{A3} }

\usepackage{preambleDJ} 

%%%%%%%%%%%% BEGIN DOCUMENT %%%%%%%%%%%%%%%
%%%%%%%%%%%% BEGIN DOCUMENT %%%%%%%%%%%%%%%
%%%%%%%%%%%% BEGIN DOCUMENT %%%%%%%%%%%%%%%
%%%%%%%%%%%% BEGIN DOCUMENT %%%%%%%%%%%%%%%
%%%%%%%%%%%% BEGIN DOCUMENT %%%%%%%%%%%%%%%

\begin{document}
\pagestyle{otherpages}
\thispagestyle{firstpage}


\vs{.6}

\section{Optimization}
\begin{enumerate}[resume]
\item A printer is asked to design a rectangular poster that will contain $1000$ square inches for the printed material, It must have $4$-inch margins at the top and bottom and $2.5$-inch margins on each side. Find the overall dimensions of the poster that minimize the total amount of paper used.
\graphicsC{poster}{2.5}
\textbf{Modeling Phase:}
\begin{enumerate}
    \item Set up an equation that represents {\em ``$1000$ square inches for the printed material"}
    \item Set up an equation that represents {\em ``total amount of paper used."}    
    \item Create the function that will be minimized {\em in terms of a single variable}.
\end{enumerate}
\textbf{Calculus Phase:}
Now that you have obtained the function to be minimized as a single variable function, use calculus to find the absolute minimum  and the corresponding dimensions of the poster.



\end{enumerate}
\newpage

\section{Optimization, again}
\enumR{
\item You are the manager of an apartment complex with all $50$ units rented at a rate $\$800$ per month.   For each increase in the rental rate of $25$ per month, one fewer apartment is rented. Maintenance costs run $\$50$ per month for each occupied unit. What is the rent that maximizes the total amount of profit?

\vspace{2mm}
\textbf{Modeling Phase:}
\begin{enumerate}
    \item Let $x$ represent the number of increases by $\$25$.
    \itemI{
    \item If $x=1$, the rental rate is $\$800+1\cdot 25=\$825$
    \item If $x=2$, the rental rate is $\$800+2\cdot 25=\$850$
    \item If $x=3$, the rental rate is $\$800+3\cdot 25=\$875$
    }
    \enum{
    \item Write down the formula for the rent per unit, $r(x)$
    \item Write down the formula for the number of occupied units $n(x)$. You will want to consider the assumption\\
     {\em ``For each increase in the rental rate of $25$ per month, one fewer apartment is rented."}. 
    }
    \item Using the results from the previous part, write down the total revenue in terms of $x$, given that
    $$
        \text{Revenue} = (\text{rent per unit}) \times (\text{number of occupied units}).
    $$
    \item Write down the formula for the total maintenance cost. You will want to consider the assumption\\ {\em ``Maintenance costs run $\$50$ per month for each occupied unit" }
    \item Write down the formula for the profit, given that
    $$
        \text{Profit} = \text{Revenue} - \text{Cost}.
    $$
\end{enumerate}

\textbf{Calculus Phase:}
Now that you have obtained the formula for the profit as a single variable function, use calculus to find the rent that maximizes the total amount of profit.
}

\newpage
\section{Team Discussion}

 As a team, talk about the similarities and differences for how you solved the optimization problem in Section 2 compared to how you solved the optimization problem in Section 3. What similar steps did you take? What was different?\\\\
Note that there is no need to record your discussion on the board and submit it to Gradescope. The point of this is to have every person on your team verbalize some/all of the steps needed to solve an optimization problem.

\newpage

\section{Probability}
\begin{enumerate}
\item Peter T, a pre-dentist student,  is a Lawn resident and was excited to participate in this year's Trick-or-Treating on the Lawn. He decided to offer Trick-or-Treaters his favorite candy bar, Snickers, as well as something, um, ``more practical"-- toothbrushes.   He put $2$ Snicker bars and $3$ toothbrushes into a grab bag and asked each Trick-or-Treater to pick two items without looking and without the ability to put any back. What's the chance a Trick-or-Treater  walked away crying because they got  $0$ Snickers (i.e $2$ toothbrushes?  What's the chance the Trick-or-Treater walked away jumping for joy because they got $2$ Snickers? What's the chance they walked away confused because they got $1$ Snickers (and thus $1$ toothbrush)?

Being a part of the famed Math $1190$ crew, you decide to create a histogram to figure out these probabilities.  
\begin{enumerate}
\item What is the sample space for pulling two items out of the bag?
\item You  let $X$ be a random variable that describes how many Snickers the Trick-or-Treater received.  What are the possible outcomes for $X$?
\item What is $P(X=0)$?  This is the chance the Trick-or-Treater gets $0$ Snickers (and thus  $2$ toothbrushes) and bursts out in tears. The next $3$ steps will help you find  $P(X=0)$.
\enum{
\item What's the chance the Trick-or-Treater will get a toothbrush on the first grab?
\item Now think about how many Snickers and how many toothbrushes remain in the bag after the first grab. What's the chance to get a toothbrush on the 2nd grab?
\item Multiply these two results together to find $P(X=0)$
}
\item What is $P(X=2)$? This is the probability the Trick-or-Treater will jump for joy because they received $2$ Snickers. The next $3$ steps will help you find  $P(X=2)$.
\enum{
\item What is the chance the Trick-or-Treater will get a Snickers on the first grab?
\item Now think about how many Snickers and how many toothbrushes remain in the bag after the first grab. What's the chance to get a Snickers on the 2nd grab?
\item Multiply these two results together to find $P(X=2 )$
}
\item What is $P(X=1)$? This is the chance that the Trick-or-Treater walked away confused because they got one Snickers and one toothbrush. {\em Hint:} There are at least 2 different approaches to come up with this probability.
\item Create a histogram that represents all of this information.\\
\begin{tikzpicture}
    \begin{axis}[
        ybar=0, % Specifies vertical bars
        ymin=0, % Ensures bars start from zero
        xtick=data, % Places x-axis ticks at data points
        xticklabels={,,}, 
                  %   xticklabels={0,1,2}, % Custom x-axis labels
                %   ytick       = {0.1,0.3,0.6},
            %yticklabels =  {$0.1$, $0.3$, $0.6$},
        enlarge x limits={0.2}, % Adjusts spacing between bars and axis limits
            y label style={anchor=south west},
            ylabel={$p(x)$},
                        xlabel = {$x$},
		x label style={anchor=west},
		 ymin = 0, ymax = .6,
        bar width=20pt, % Adjusts the width of the bars
                    width = 2.5in, height=1.4in,
       % legend style={at={(0.5,-0.15)},anchor=north,legend columns=-1} % Legend positioning
    ]
  \addplot[fill=blue!40] coordinates {(1,0)  (2,0)  (3,0) };
  %   \addplot[fill=blue!40] coordinates { (1,0.3)  (2,0.6)  (3,0.1) };
    \end{axis}
\end{tikzpicture}
\end{enumerate}
\end{enumerate}

\newpage


\section{Probability, part deaux}
\begin{enumerate}[resume]
    \item Suppose two standard dice are rolled and the result is recorded.
\begin{enumerate}
    \item[a)] Let $X$ be the difference of first die minus the second die. 
    \enum{
    \item Construct a table of all possible combinations of rolls of dice and the resulting difference similar to the one we used in class.
    \item Find the probability for all possible outcomes for $X$
    \item Sketch a histogram
    }
    \end{enumerate}
    \mpageT{.5}{.5}[\hs{.2}]{
    \begin{tikzpicture}
    \begin{axis}[
        ybar, % Specifies vertical bars
        bar shift=0,
        ymin=0, % Ensures bars start from zero
        xtick=data, % Places x-axis ticks at data points
                     xticklabels={,,,,,,,,,}, % Custom x-axis labels
                     ytick=\empty,
%                    ytick       = {1/36, 2/36, 3/36, 4/36, 5/36, 6/36},
%            yticklabels =  {$1/36$,$2/36$,$3/36$, $4/36$,$5/36$,$6/36$},
        enlarge x limits={0.2}, % Adjusts spacing between bars and axis limits
            y label style={anchor=south west},
            ylabel={$p(x)$},
                        xlabel = {$x$},
		x label style={anchor=west},
		 ymin = 0, ymax = .17,
        bar width=17pt, % Adjusts the width of the bars
                    width = 4in, height=1.4in,
                    %enlarge x limits={1.2}
       % legend style={at={(0.5,-0.15)},anchor=north,legend columns=-1} % Legend positioning
    ]
 \addplot[fill=blue!40] coordinates {(1,0)  (2,0)  (3,0) (4,0) (5,0) (6,0) (7,0) (8,0) (9,0) (10,0) (11,0)};
 %\addplot[fill=blue!40] coordinates {(1,0.027)  (2,0.055)  (3,0.083) (4,0.11) (5,0.139) (6,0.167) (7,0.139) (8,0.11) (9,0.083) (10,0.055) (11,0.027)};
    \end{axis}
\end{tikzpicture}
    }{
    \begin{enumerate}
    \item[b)] Compute $P(X \geq 2)$.
    \item[c)] Compute $P(X \leq -2)$.
    \item[d)] Compute $P(-1\leq X \leq 1)$.
    \item[e)] Compute $P(X \text{ is even})$.
    \end{enumerate}
    }




\end{enumerate}


\end{document}